\subsection{One deleted addition and exponential cover gain}
A single deleted addition imposes a stronger restriction than the small-profile clause counts alone suggest.

\begin{theorem}[A single deleted addition]\label{thm:purecut}
Let $G$ be a finite simple graph without a four-cycle, and let a positive variable-input binary circuit compute $P\in\Sup(G)$ with nonempty quadratic support $H$. Suppose exactly one of the gates deleted by minimum-degree extraction and pruning is an addition. Then, for every vertex cover $Z$ of $H$, the induced graph $G[V(G)\setminus Z]$ is a matching with possible isolated vertices. Moreover,
\[
 D\ge 2^d-1,\qquad d=\tau(G)-\tau(H).
\]
If $G$ has girth at least five, the missing-edge graph $E(G)\setminus E(H)$ itself is a matching.
\end{theorem}
\begin{proof}
If no edge is missing, the induced graph outside any $H$-cover is empty and the conclusions are immediate. Otherwise apply full-carrier erasure. The sole deleted addition cannot disappear, since its disappearance would leave no route for a missing witness to enter the selected quadratic output. It is the unique visible cut, and every useful private product is ancestral to its high operand $g$. Old quadratic carriers are unavailable as off-carrier inputs. Thus $g$ is a multiplication DAG over old linear boundary values. Booleanly it is a conjunction of their positive clauses.

The independent-symbol quadratic frame is affine in this one high signal. If a degree-one symbol occurred squared, substituting its nonzero old linear low part would give a square in the quadratic output. A degree-two symbol has only scalar context. Hence, after setting any $H$-cover $Z$ to zero, the quadratic part disappears and the output is $gL$, where $L$ is zero, one, or a positive linear clause. This is a positive CNF for the edge function of $G[T]$, where $T=V(G)\setminus Z$.

Suppose $u,w,v\in T$ form a two-edge path, with $u\ne v$. No old boundary clause can contain both $u$ and $v$. A clause containing both variables uses a retained linear-addition gate. That gate reaches a quadratic product, so some opposite-factor variable $h$ is an $H$-neighbor of both. Because $Z$ covers $H$ and $u,v\in T$, the vertex $h$ lies in $Z$ and is distinct from $u,w,v$. The four edges $uw,wv,vh,hu$ would form a four-cycle. The same argument applies to a nonconstant direct context $L$: for its nonzero old low part $A$, the positive product $AL$ occurs in the quadratic frame, so a variable in $A$ would give the same common old neighbor.

Extend $\{w\}$ to a maximal independent set $I$ of $G[T]$. Its complement $U=T\setminus I$ is an inclusion-minimal vertex cover containing $u,v$. Since the CNF $gL$ is false on $I$, some factor clause has support contained in $U$. Every clause of a positive CNF representing an edge function is itself a vertex cover, because it must be true on every edge assignment. Minimality of $U$ forces that clause's support to equal $U$. It would contain $u,v$, although no old factor or the single context can do so. This contradiction proves that $G[T]$ has maximum degree at most one.

Choose $Z$ minimum and let $m$ be the number of edges of the resulting matching. Adding one endpoint of each to $Z$ covers $G$, so $m\ge d$. If $r$ private products remain after erasure, their connected single-output multiplication DAG has at most $r+1$ distinct old boundary sources. Indeed, its $r$ binary nodes have $2r$ input occurrences, at least $r-1$ of which join product nodes; at most $r+1$ boundary occurrences remain. Reuse of a source reduces the number of distinct factors. The context adds at most one clause, so $gL$ has at most $r+2$ clauses.

A matching of $m$ edges has $2^m$ minimal vertex covers. Each is the complement of a maximal independent assignment and must occur as a clause in any positive CNF representing the matching, by the preceding minimality argument. Therefore
\[
 2^d\le2^m\le r+2\le D+1,
\]
since erasure cannot increase the original $D-1$ deleted products. This proves the displayed lower bound without any quadratic-optimality hypothesis.

Finally suppose $G$ is also triangle-free. If two missing edges $uw,wv$ met, there would be no edge $uv$. The three vertices would contain no $H$-edge, so $V(G)\setminus\{u,w,v\}$ would be an $H$-cover. The first conclusion would contradict the two-edge path on its complement. Thus the missing edges are pairwise disjoint.
\end{proof}

In particular, cover gain three with five or six deletions requires at least two deleted additions, whatever the quadratic excess and the shape of the missing edges. This applies to every remaining cover-gain-three alternative in Theorem~\ref{thm:four-summary}.
